Kita meninjau pemetaan diferensiabel
\maabbeledisp {\varphi} {\R^3} {\R
} {(x,y,z)} { x^2+y^4+z^6-1
} {.}
Himpunan $M$ adalah serat $\varphi$ di atas $0$. Berlaku

\mathdisp {\left(D\varphi\right)_{(x,y,z)}  = (2x,4y^3,6z^5)} { . }

Turunan ini sama dengan $(0,0,0)$ hanya di
\mathl{(x,y,z)=(0,0,0)}{}, dan titik ini bukan anggota $M$, sehingga $\varphi$ reguler di setiap titik $M$. Karena itu, menurut teorema pemetaan implisit, himpunan tersebut merupakan sebuah manifold berdimensi dua.

Sebagai serat dari sebuah pemetaan kontinu, $M$ merupakan sebuah himpunan bagian tertutup dari $\R^3$. Selain itu, $M$ terbatas. Memang, untuk
\mathl{(x,y,z) \in M}{} berlaku
\mathl{\betrag {  x  } , \betrag {  y  } ,\betrag {  z  } \leq 1}{,} sebab jika tidak,
\mathl{x^2+y^4+z^6 >1}{}. Hal ini mengimplikasikan kekompakan.

 [[Kategorie:Latexseite]]
